\begin{abstract}
Pre-deployment estimates could screen expensive LLM inference experiments, but
an unverified estimate is not measurement evidence.  We present \system{}, a
lightweight, evidence-gated sandbox that combines an interpretable analytical
projector, a sparse correction learned from measured workloads, and an
independent validation boundary.  Every scenario, prediction, telemetry trace,
and measurement is split-labeled and content-addressed; a prediction must be
frozen before the corresponding hardware run.  On one H100 80GB serving
Qwen2.5-7B with vLLM, we use one three-repeat anchor and six development
workloads to calibrate workload transfer across context length and concurrency.
We then freeze eight unseen workloads and collect three independent repeats per
point.  Relative to the analytical projector, the calibrated sandbox reduces
blind energy MAPE from \BaseEnergyMAPE{} to \CalibratedEnergyMAPE{}, preserves
energy ordering with Spearman $\rho=\EnergySpearman{}$ and
\EnergyPairwisePct{} pairwise accuracy, and identifies the measured
lowest-energy workload.  Throughput and TPOT MAPE are
\CalibratedThroughputMAPE{} and \CalibratedTPOTMAPE{}; TTFT remains harder at
\CalibratedTTFTMAPE{}, exposing a low-concurrency boundary.  The result supports
cheap workload screening within the measured envelope, not multi-GPU or
cross-hardware simulation, and demonstrates why explicit evidence gates are
necessary before an optimizer or agent consumes sandbox outputs.
\end{abstract}

\begin{IEEEkeywords}
LLM inference, energy efficiency, workload modeling, multi-fidelity evidence,
GPU measurement
\end{IEEEkeywords}
